\section{Singular Fourier tests for the archimedean vacuum}
\label{sec:sw-singular-vacuum}

Suppose \(0\le m<n\). In the definite Siegel--Weil calculation of
Section~\ref{sec:sw-boundary}, the local vacuum at the special
parameter has Fourier frequencies of rank at most \(m\). We prove
this assertion after arbitrary left and right translations and
after averaging by a compact group. The singularities here concern
the rank of the Fourier frequency.

Use \(J=\left(\begin{smallmatrix}0&I\\-I&0\end{smallmatrix}\right)\)
and \(G=\{g\in\operatorname{GL}_{2n}(\mathbb C):g^*Jg=J\}\).
Write \(A_g,B_g,C_g,D_g\) for its four \(n\)-by-\(n\) blocks, and set
\[
 U_g=iC_g+D_g,\quad V_g=iA_g+B_g,\quad
 n(S)=\begin{pmatrix}I&S\\0&I\end{pmatrix},\qquad
 f(g)=\det(g)^e\det(U_g)^{-m},\qquad e\in\mathbb Z.
\]
Here \(f\) is \texttt{fz} at parameter zero, the local parameter
measured relative to the distinguished global value. Identify
\(s\in\mathbb R^{n\times n}\) with the Hermitian matrix
\[
 S(s)_{jk}=\tfrac12(s_{jk}+s_{kj})+
                   \tfrac i2(s_{jk}-s_{kj}).
\]
This real-linear bijection fixes the coordinate measure \(ds\).
Call \(\psi\in C_c(\mathbb R^{n\times n},\mathbb C)\) a singular
test if for each Hermitian \(T\) with \(\det T=0\),
\[
       \int \psi(s)\exp(2\pi i\operatorname{tr}(T S(s)))\,ds=0 .
\]
Then for each \(y,b\in G\),
\begin{equation}
           \int f\bigl(y n(S(s)) b\bigr)\psi(s)\,ds=0.
 \label{eq:sw-singular-vacuum}
\end{equation}

\paragraph{A matrix Gaussian with singular frequencies.}
For a positive definite Hermitian \(P\), a Hermitian \(H\), and
complex Lebesgue measure in the \(nm\) entries of an
\(n\)-by-\(m\) matrix \(Z\), one has
\begin{equation}
 \det(P-iH)^{-m}
 =\pi^{-nm}\int_{\mathbb C^{n\times m}}
   e^{-\operatorname{tr}(Z^*PZ)}e^{i\operatorname{tr}(Z^*HZ)}\,dZ.
 \label{eq:sw-singular-gaussian}
\end{equation}
Indeed, put \(Z=P^{-1/2}W\). The real Jacobian for each column is
\(\det P^{-1}\). Diagonalize the Hermitian
\(P^{-1/2}HP^{-1/2}\) by a unitary change preserving both measure
and the norm. Each eigenvalue \(\lambda\) gives a factor
\[
 \int_{\mathbb C} e^{-(1-i\lambda)|z|^2}\,d\Re z\,d\Im z
       =2\pi\int_0^\infty r e^{-(1-i\lambda)r^2}\,dr
       =\frac{\pi}{1-i\lambda}
\]
in each column. Multiplication of these factors proves the
formula, including the normalization. All changes and factorizations
are justified by the absolute Gaussian majorant. For \(m=0\) the
integral and the determinant power are both one.

Fix a Hermitian \(H_0\). Insert \(H=H_0+S(s)\) in
\eqref{eq:sw-singular-gaussian} and integrate against \(\psi(s)\).
The absolute integral is bounded by
\(\|\psi\|_1\int e^{-\operatorname{tr}(Z^*PZ)}\,dZ<\infty\).
For fixed \(Z\), the \(S(s)\)-dependent factor becomes
\(\exp(2\pi i\operatorname{tr}(T S(s)))\) with
\(T=ZZ^*/(2\pi)\). It is Hermitian and has rank at most \(m<n\).
The definition of a singular test therefore annihilates the inner
\(s\)-integral, and Fubini gives
\begin{equation}
      \int\psi(s)\det(P-i(H_0+S(s)))^{-m}\,ds=0.
 \label{eq:sw-singular-hermitian}
\end{equation}
In particular the zero frequency is included. Transposing the
rectangular coordinate array gives the row convention of the
source, with the same trace pairing and determinant.

\paragraph{Factorization on the open matrix cell.}
Let \(y,b\in G\) and first assume \(C=C_y\) invertible. Write
\(D=D_y\), \(U=U_b\), \(V=V_b\). Since
\(\binom{V}{U}=b\binom{iI}{I}\), preservation of \(J\) gives
\[
          U^*V-V^*U=2iI .
\]
If \(Uv=0\), evaluating the identity against \(v\) gives
\(2i\|v\|^2=0\), so \(U\) is invertible. This argument holds for
every \(b\in G\). Also \(yJy^*=J\), since \(J^{-1}=-J\) and
\(y^*Jy=J\). Its bottom-right block gives \(CD^*=DC^*\).
It follows that \(C^{-1}D=(C^{-1}D)^*\).

Put \(L=VU^{-1}\). Multiplying the preceding identity for \(U,V\)
by \((U^{-1})^*\) and \(U^{-1}\) shows that
\[
 L-L^*=2iP,\qquad P=(U^{-1})^*U^{-1}>0,\qquad
 H_0=\tfrac12(L+L^*)+C^{-1}D=H_0^* .
\]
For any Hermitian \(S\), direct block multiplication gives
\[
 \begin{split}
 U_{y n(S)b}
     &=CV+(D+CS)U\\
     &=C(L+C^{-1}D+S)U\\
     &=C\bigl(i(P-i(H_0+S))\bigr)U .
 \end{split}
\]
Thus
\(\det U_{y n(S)b}=\kappa\det(P-i(H_0+S))\), where
\(\kappa=i^n\det C\det U\ne0\).
Since \(\det n(S)=1\),
\[
 f(y n(S)b)
   =(\det y\det b)^e\,\kappa^{-m}\det(P-i(H_0+S))^{-m}.
\]
Its integral vanishes by \eqref{eq:sw-singular-hermitian}.

\paragraph{The remaining matrices and the compact average.}
For arbitrary \(y\in G\), multiply on the right by
\(\bar n(t)=\left(\begin{smallmatrix}I&0\\tI&I\end{smallmatrix}\right)\)
with \(t\) real. This belongs to \(G\), and the new lower-left block
is \(C_y+tD_y\). The polynomial \(p(T)=\det(C_y+TD_y)\) is nonzero,
since \(p(-i)=(-i)^n\det(iC_y+D_y)\ne0\). Its real roots form a
finite set. Thus for all sufficiently small nonzero real \(t\)
the result on the open cell applies to \(y\bar n(t)\).
The determinants defining \(f\) are everywhere nonzero on \(G\),
so \((t,s)\mapsto f(y\bar n(t)n(S(s))b)\) is continuous. It is
bounded on \([-1,1]\times\operatorname{supp}\psi\).
Dominated convergence as \(t\to0\) proves
\eqref{eq:sw-singular-vacuum} for the original \(y\).

Let \(K\) be a compact group with a continuous homomorphism
\(\iota\) onto \(G\cap U(2n)\), let \(\mu\) be its Haar measure,
and let \(\rho\) be a continuous complex weight on \(K\).
For \(a,b,x\in G\), the associated averaged test has value
\[
 \int_K \rho(k)\left(\int
       f(x\iota(k)a n(S(s)) b)\,\psi(s)\,ds\right)d\mu(k)=0.
\]
Indeed the inner integral vanishes at each \(k\) by
\eqref{eq:sw-singular-vacuum}, with \(y=x\iota(k)a\).
The product of \(K\) and the support of \(\psi\) is compact,
so the double integral is absolutely convergent. This covers
the allowed weights of the local vacuum projection.
The argument evaluates the section at its distinguished parameter.
Derivatives in the parameter, and their use in controlling poles
of the continued Eisenstein series, require an additional argument.

\paragraph{Source correspondence.}
The tagged source defines \texttt{IsSingTest} using the displayed
coordinate pairing in
\srcid{SiegelWeilProof.HolomorphicContinuation.IsSingTest}.
Within its \texttt{CayleyCosetSums} directory, the rank bound and
Gaussian integral appear in \texttt{TraceTransposeMul} and
\texttt{InnerGen.core2}; the block identities and factorization
are \texttt{SecondCoordMulInner} and \texttt{Imports.factor}.
The extension from the open cell is
\texttt{VanishInv.vanish\_inv} and
\texttt{EventuallyUnit.vanish\_all}. Its compact-average conclusion
is \texttt{SingularTestsKillVacuum.lt3}.
\footnote{\href{https://github.com/ChristianKudeba/alpoge-rh-verification/tree/alpoge-source-7cbf7a0c/src/Development/WeilBound/RepresentationTheoryGL/SiegelWeilProof/HolomorphicContinuation/CayleyCosetSums}{Tagged local proof files}.
The global unitary identity in the anisotropic range is
S.~Yamana, \emph{The Siegel--Weil formula for unitary groups},
Pacific J. Math. 264 (2013),
\href{https://msp.org/pjm/2013/264-1/pjm-v264-n1-p11-p.pdf}{Theorem~2, p.~237}.}
